% !TeX root = ../../金融经济学_王江_讲义及习题.tex
% 知识点：references/金融经济学/金融经济学-王江_261013.pdf，第 2 章
% 补充：references/金融经济学/金融经济学_第2章 基本框架.pptx（第 1–44 页，含 WMF 图片公式）
% 习题：references/金融经济学/史上最全最完整的金融经济学答案(王江版).pdf，页 2–6（印刷页 1–5）

\lecture{基本框架}

\ifdefined\WJHandoutKnowledge

\begin{tcolorbox}[knowledge,title=2.1 经济环境]
本章先描述经济的三个基本要素：环境、参与者和证券市场，再讨论个人选择、市场均衡与配置效率。环境的两个关键维度是\textbf{时间与风险}：0 期交易和消费，1 期自然状态实现，再按该状态兑现支付。不同日期、不同状态的消费是不同的或有商品。
\begin{enumerate}

\item 形式上可以用一棵状态树来表示上面描述的经济环境。用 $\omega$ 标识状态，假设总共有 $\Omega$ 个状态，那么表示所有状态可能性的集合 $\Omega$ 称为\textbf{状态空间}。记状态发生的概率为 $p(\omega)$，概率测度记为 $P=\{p_{\omega}, \omega\in\Omega\}$，其中 $p_\omega\ge0$、$\sum_\omega p_\omega=1$。同一符号 $\Omega$ 在教材中兼指状态集合与状态数，须由上下文区分。

\item 假设经济中只有一种不可储存的商品，且本模型没有参与者可选择的生产技术。因此不能直接把今天的商品储存到明天；证券交易转移的是对各日期、各状态资源的要求权，不改变给定状态下的总资源。
\item 用事件树表示未来状态及其概率：
\begin{center}
\WJEventTree[left={$\Omega$},stem]{
    \omega=1 \text{，概率为 }p_1,
    \omega=2 \text{，概率为 }p_2,
    \vdots,
    \omega=\Omega \text{，概率为 }p_{\Omega}
}
\end{center}
\item \textbf{例 2.1（Lucas 树经济）：}一棵桃树今年产出 $100$，明年好天气产出 $200$、坏天气产出 $50$，两状态等概率，此后不再产出。状态编号与产出数量应分开：$\omega=a,b$ 是标签，$Y_1(a)=200$、$Y_1(b)=50$ 才是支付。
\begin{center}
\WJEventTree[left={$100$},stem]{200\text{，状态 }a\text{，概率 }1/2,50\text{，状态 }b\text{，概率 }1/2}
\end{center}
\end{enumerate}
\end{tcolorbox}


\begin{tcolorbox}[knowledge,title=2.2 经济参与者]
    把经济参与者描述为：为自己切身利益而行动的个体。一个参与者的经济特征包括：
    \begin{enumerate}
        \item 他控制的经济资源（任何有经济价值，可以用于满足经济需求的东西，可以是有形 or 无形的）。参与者的经济资源包括
        \begin{itemize}
            \item 所占有的实物商品/资本品
            \item 信息
            \item 生产技术
        \end{itemize}
        \item 他的经济需求。
    \end{enumerate}

\tcbline

参与者的经济资源：
\begin{enumerate}
    \item 禀赋：即参与者初始占有的资源，是模型开始时归他所有的商品或资本品（不要求来自出生时）。给定两个时期，且第二时期有多个状态，那么他的禀赋可表示为：
    \[
    e_k=\begin{bmatrix}e_{k,0} \\ e_{k,1}\end{bmatrix}
    =\begin{bmatrix}e_{k,0} \\ \begin{bmatrix}e_{k,11} \\ \vdots \\ e_{k,1\Omega}\end{bmatrix}
    \end{bmatrix}
    \quad k=1,\cdots,K
    \]

    在数学上，可以把参与者的禀赋看成是 $1+\Omega$ 维实空间 $\mathrm{R}^{1+\Omega}$ 中的一个元素，且假设参与者的禀赋非负，那么禀赋 $e_k \in \mathrm{R}^{1+\Omega}_{+}$。
\begin{tcolorbox}[warning,title=向量的“大于”]
    对于向量 $a=\begin{bmatrix}a_1 \\ a_2 \\ \vdots \\ a_n\end{bmatrix}$，有以下定义：
    \begin{enumerate}
        \item $a \ge 0$，即 a 的所有分量都非负，即 $a_i \ge 0, i=1,\cdots,n$。
        \item $a > 0$，即 a 的所有分量都非负，且至少有一个分量严格大于 0，即 $\exists i \in \{1,\cdots,n\}$ 使得 $a_i > 0$。
        \item $a \gg 0$，即 a 的所有分量都严格大于零，即 $a_i > 0, i \in \{1,\cdots,n\}$。
    \end{enumerate}
\end{tcolorbox}
    \item 信息：信息具备经济价值。参与者的信息可以分为两种类型：公开信息和私有信息。如果参与者仍然拥有私有信息，信息就是不对称的。信息不对称会导致很多问题，例如道德风险和逆向选择。但在本书中，将假设经济者对于经济，特别是其未来状态发生的可能性，拥有相同的信息。
    \item 生产技术：指可以把商品转化为更多商品，或者把现在的商品转化为未来的商品。在本书中，假设所有参与者不拥有任何生产技术。
\end{enumerate}

\tcbline

参与者的经济需求：
参与者的经济需求基于他的不同时期和未来不同状态下的消费。对于参与者 k，记他在第零期的消费是$c_{k,0}$，第一期的消费是 $c_{k,\omega}$，其中 $\omega \in \Omega$。消费的完整表述是：
\[
c_k \equiv \begin{bmatrix}c_{k,0} \\ c_{k,1}\end{bmatrix}=
\begin{bmatrix}c_{k,0} \\ \begin{bmatrix}c_{k,11} \\ \vdots \\ c_{k,1\omega} \\ \vdots \\ c_{k,1\Omega}\end{bmatrix}
\end{bmatrix} \quad k=1,\cdots,K
\]


\begin{enumerate}
\item 消费集：消费者可能的消费选择叫做消费计划 $c=\begin{bmatrix}
    c_0 \\ c_1
\end{bmatrix}$。所有消费计划的集合叫做消费集 $C$。和禀赋一样，消费计划属于 $\mathrm{R}^{1+\Omega}$ 中的一个元素，且假设消费计划非负，那么消费计划 $c \in \mathrm{R}^{1+\Omega}_{+}$，消费集就是 $C=\mathrm{R}^{1+\Omega}_{+}$。

\begin{tcolorbox}[knowledge]
\textbf{定义 2.1} 设 A 是 $R_n$ 的一个集合，对于 $\forall a,b \in A$ 和 $\alpha \in [0,1]$，如果 $a\alpha+(1-\alpha)b$ 也在A 中，则称集合 A 是凸的（convex）。

\textbf{定义 2.2} 对于 $\mathrm R^n$ 的集合 $A$，如果其中任意\textbf{收敛}序列 $a_i\to a$ 的极限也在 $A$ 中，则称集合 A 是闭的（closed）。

\textbf{假设 1} 消费集 $C=\mathrm{R}^{1+\Omega}_{+}$ 是 $\mathrm{R}^{1+\Omega}$ 中的一个闭凸子集。

\textbf{直接验证：}若 $a,b\ge0$、$\alpha\in[0,1]$，则每个分量 $\alpha a_i+(1-\alpha)b_i\ge0$，故 $C$ 凸。若 $c^{(n)}\in C$ 且 $c^{(n)}\to c$，则 $c_i=\lim_{n\to\infty}c_i^{(n)}\ge0$，故 $c\in C$，所以 $C$ 闭。“闭”不要求集合中每个序列都有极限，只要求已有极限不落到集合外。
\end{tcolorbox}

\item 偏好：偏好是参与者对所有可能性计划的排序这样的排序定义了他的经济需求。

\begin{tcolorbox}[knowledge]
\textbf{定义 2.3} 偏好是C 上的一个二元关系，表示为 $\succeq$，且满足以下三个条件：
\begin{itemize}
    \item 完备性：对于 $\forall a,b \in C$，有 $a\succeq b$ 或 $b\succeq a$。
    \item 传递性：对于 $\forall a,b,c \in C$，如果 $a\succeq b$ 且 $b\succeq c$，则 $a\succeq c$。
    \item 自反性：对于 $\forall a \in C$，有 $a\succeq a$。
\end{itemize}
\end{tcolorbox}

记 $a\succ b$ 表示 $a\succeq b$ 且不是 $b\succeq a$（严格偏好）；$a\sim b$ 表示二者相互弱偏好（无差异）。这是对整段消费计划的排序，不能与向量逐分量比较 $a>b$ 混为一谈。

\item 偏好的基本假设：\textcolor{red}{\textbf{参与者喜多厌少}：如果两个消费计划 a 和 b 之间，在所有时期和状态下，a 都不比 b 提供的消费少，且在有些状态下提供更多的消费，那么 a 优于 b。}
    \begin{tcolorbox}[knowledge]
        \textcolor{red}{\textbf{公理 1 不满足性：}} 如果 $a>b$，那么 $a\succ b$（严格偏好）。

        \textbf{公理 2 连续性：} $\forall c \in C$，集合 $\{a\in C: a\succeq c\}$ 和 $\{b\in C: b\preceq c\}$ 都是闭集。

        \textbf{公理 3 凸性：} $\forall a,b \in C$ 以及 $\alpha \in [0,1]$，如果 $a \succeq b$ ，那么有 $\alpha a+(1-\alpha)b\succeq b$ 。
    \end{tcolorbox}

    \item 效用函数：从消费计划到实数的映射叫做效用函数。
    
    \begin{tcolorbox}[knowledge]
        \textbf{定义 2.4} 对应于偏好关系的效用函数 U 是从 C 到 R 的函数，使得 $\forall a,b \in C$，当且仅当 $a\succeq b$ 时，有 $U(a) \ge U(b)$。
    \end{tcolorbox}

    \begin{tcolorbox}[knowledge]
        \textbf{定理 2.1} 对于一个在闭凸消费集 C 上，有定义 2.3 定义的（完备性、传递性和自反性）、满足公理 2 的偏好（连续性），存在一个连续的效用函数 U 使得 $\forall a,b \in C, a\succeq b$，当且仅当 $U(a) \ge U(b)$。
    \end{tcolorbox}

    \begin{tcolorbox}[example]
        \begin{question}
            \textbf{例 2.3}\quad 在例 2.1 中描述的 Lucas 树经济中，消费空间是 $\mathrm{R}^{3}_{+}$，表示在今年和明年的两个可能状态下的消费；状态 $a$ 对应于好天气，而状态 $b$ 则对应于坏天气。假设经济中的某个参与者有如下的效用函数：
            \[
                U(c_0,c_{1a},c_{1b})=\log c_0+\tfrac12(\log c_{1a}+\log c_{1b})
            \]
            在如下三个消费计划中，他将会选择哪一个？
            \begin{center}
            \begin{tabular}{@{}c@{\quad}c@{\quad}c@{}}
                \textbf{计划 A} & \textbf{计划 B} & \textbf{计划 C} \\[0.35em]
                \WJEventTree[left={$1$},stem]{1,1} &
                \WJEventTree[left={$\tfrac12$},stem]{\tfrac32,\tfrac32} &
                \WJEventTree[left={$1$},stem]{\tfrac32,\tfrac12}
            \end{tabular}
            \end{center}
        \end{question}

        \begin{solution}
            对于 A、B 和 C，效用函数给出的相应的效用值分别为 $U_A=0$、$U_B=\log(3/4)\approx-0.288$、$U_C=\tfrac12\log(3/4)\approx-0.144$。显然，参与者对三个消费计划的排序是 $A>C>B$，因此他会选择 A。
        \end{solution}

        \begin{question}
            \textbf{例 2.4}\quad 让我们考虑同一经济中的另外一个参与者，他的效用函数是
            \[
                V(c_0,c_{1a},c_{1b})=c_0\sqrt{c_{1a}c_{1b}}
            \]
            他对上面三个消费计划的排序是什么？
        \end{question}
        \begin{solution}
            \[
                \begin{aligned}
                    V_A &= 1\cdot\sqrt{1\cdot 1}=1,\\
                    V_B &= \tfrac12\cdot\sqrt{\tfrac32\cdot\tfrac32}=\tfrac34,\\
                    V_C &= 1\cdot\sqrt{\tfrac32\cdot\tfrac12}=\tfrac{\sqrt3}{2}.
                \end{aligned}
            \]
            因此 $V_A>V_C>V_B$，参与者对三个消费计划的排序仍是 $A>C>B$。
        \end{solution}
    \end{tcolorbox}
\end{enumerate}

上面的例子说明了如下的结论：从定理 2.1 中得到的效用函数只是\textbf{序数的}（ordinal），即它只给出了消费计划的排序。效用函数的任意\textbf{严格递增}变换不会改变这种排序：若 $V=f(U)$、$f$ 严格递增，则 $V(a)\ge V(b)\iff U(a)\ge U(b)$。本例在正消费域内恰有
\[
 e^{U(c)}=e^{\log c_0+\frac12\log c_{1a}+\frac12\log c_{1b}}
 =c_0\sqrt{c_{1a}c_{1b}}=V(c),
\]
所以两种公式给出相同偏好。效用差额或倍数没有直接福利数量含义，也不能直接进行人际比较。对数效用只在正消费域定义，不能据此宣称它在包含零消费的整个 $\mathrm R^{1+\Omega}_+$ 上连续。

利用效用函数作为偏好的一种表述形式，我们可以把偏好的三个基本假设重新表示如下：
\begin{enumerate}
    \item \textbf{不满足性：}如果 $a>b$，那么 $U(a)>U(b)$；相应有 $a\ge b\Rightarrow U(a)\ge U(b)$。
    \item \textbf{连续性：}$\lim_{a_n \to a}U(a_n)=U(a)$。
    \item \textbf{凸性：}对任意 $a,b$ 和 $\alpha\in[0,1]$，
    \[
      U(\alpha a+(1-\alpha)b)\ge\min\{U(a),U(b)\}.
    \]
    这是\textbf{拟凹性}，即上等值集为凸集，并非 $U$ 作为函数的凸性。若选取凹效用表示，则更强的 Jensen 不等式 $U(\alpha a+(1-\alpha)b)\ge\alpha U(a)+(1-\alpha)U(b)$ 足以推出偏好凸。严格递增变换保留排序与拟凹性，却未必保留函数的凹性；不能把任一序数表示的曲率直接当作风险厌恶。
\end{enumerate}

本章习题里反复出现的是下面几种效用函数$u(x)$的形式。两状态等概率时，跨期效用写成
\[
    U(c)=u(c_0)+\tfrac12\bigl(u(c_{1a})+u(c_{1b})\bigr).
\]

表中只记函数形状，以及它如何对应不满足性和凸性。绝对风险厌恶、相对风险厌恶以及 CARA、CRRA 的比较留到第 6、7 章。

{\small
\renewcommand{\arraystretch}{1.25}
\noindent\begin{tabularx}{\linewidth}{@{}llX@{}}
\toprule
名称 & 效用函数 $u(x)$ & 要点 \\
\midrule
对数 & $\log x$ & $x>0$。严格递增、严格凹。幂效用去掉与消费无关的常数后，在 $\gamma\to1$ 时趋于对数。 \\
幂 & $\dfrac{x^{1-\gamma}}{1-\gamma}$ & $\gamma>0$、$\gamma\neq 1$，$x>0$。严格递增、严格凹。平方根效用即 $\gamma=\tfrac12$。常相对风险厌恶（CRRA）。 \\
负指数 & $-e^{-ax}$ & $a>0$。严格递增、严格凹。常绝对风险厌恶（CARA）。 \\
二次 & $x-\tfrac12 a x^{2}$ & $a=0$ 时线性，处处递增；$a>0$ 时仅当 $x<1/a$ 边际效用为正，不满足性只在此范围内成立。 \\
\bottomrule
\end{tabularx}}

下图只示意形状，纵轴各自平移和缩放。幂效用取 $\gamma=2$，负指数与二次取 $a=1$。二次图中的虚线是边际效用降为零的位置 $x=1/a$。
\begin{center}
\setlength{\tabcolsep}{2pt}
\footnotesize
\begin{tabular}{@{}cccc@{}}
\begin{tikzpicture}[baseline=(current bounding box.center), x=1cm, y=1cm]
    \path (-0.28,-0.42) rectangle (2.62,1.72);
    \draw[-{Stealth[length=1.5mm]}, thin] (0.22,0) -- (2.32,0) node[below left, inner sep=1pt] {$x$};
    \draw[-{Stealth[length=1.5mm]}, thin] (0.22,0) -- (0.22,1.48) node[above left, inner sep=1pt] {$u$};
    \draw[blue!65!black, thick, domain=0.32:2.15, samples=50, smooth] plot (\x,{0.95+0.62*ln(\x)});
\end{tikzpicture}
&
\begin{tikzpicture}[baseline=(current bounding box.center), x=1cm, y=1cm]
    \path (-0.28,-0.42) rectangle (2.62,1.72);
    \draw[-{Stealth[length=1.5mm]}, thin] (0.22,0) -- (2.32,0) node[below left, inner sep=1pt] {$x$};
    \draw[-{Stealth[length=1.5mm]}, thin] (0.22,0) -- (0.22,1.48) node[above left, inner sep=1pt] {$u$};
    \draw[blue!65!black, thick, domain=0.34:2.15, samples=50, smooth] plot (\x,{1.42-0.40/\x});
\end{tikzpicture}
&
\begin{tikzpicture}[baseline=(current bounding box.center), x=1cm, y=1cm]
    \path (-0.28,-0.42) rectangle (2.62,1.72);
    \draw[-{Stealth[length=1.5mm]}, thin] (0,0) -- (2.32,0) node[below left, inner sep=1pt] {$x$};
    \draw[-{Stealth[length=1.5mm]}, thin] (0,0) -- (0,1.48) node[above left, inner sep=1pt] {$u$};
    \draw[blue!65!black, thick, domain=0:2.15, samples=40, smooth] plot (\x,{1.32-1.12*exp(-\x)});
\end{tikzpicture}
&
\begin{tikzpicture}[baseline=(current bounding box.center), x=1cm, y=1cm]
    \path (-0.28,-0.42) rectangle (2.62,1.72);
    \draw[-{Stealth[length=1.5mm]}, thin] (0,0) -- (2.32,0) node[below left, inner sep=1pt] {$x$};
    \draw[-{Stealth[length=1.5mm]}, thin] (0,0) -- (0,1.48) node[above left, inner sep=1pt] {$u$};
    \draw[blue!65!black, thick, domain=0:2, samples=40, smooth] plot ({1.05*\x},{2.4*(\x-0.5*\x*\x)});
    \draw[dashed, black!45] (1.05,0) -- (1.05,1.2);
    \node[below, inner sep=1pt] at (1.05,0) {$1/a$};
\end{tikzpicture}
\\[0.15em]
对数 $\log x$ & 幂 $\dfrac{x^{1-\gamma}}{1-\gamma}$ & 负指数 $-e^{-ax}$ & 二次 $x-\tfrac12 ax^{2}$
\end{tabular}
\end{center}

\end{tcolorbox}



\begin{tcolorbox}[knowledge,title=2.3 证券市场]
资源的配置是通过金融市场中的交易来完成的，市场允许参与者在不同时期和状态上配置资源。

\tcbline

\begin{enumerate}
    \item \textbf{证券及其支付：}在我们所要构建的简单框架中，金融市场由一组\textbf{证券}（security）构成。一只证券即是一份\textbf{金融要求权}（financial claim），它在 1 期给其所有者带来\textbf{支付}（payoff）；支付的数量一般依赖于当时的经济状态。令 $x_\omega$ 为证券在状态 $\omega$ 时的支付，其中 $\omega\in\Omega$，则一只证券由它在各个可能状态下的支付 $\{x_\omega,\,\omega=1,\dots,\Omega\}$ 来定义。形式上可用一棵\textbf{支付树}表示：
    \begin{center}
    \WJEventTree[left={$0$},stem]{x_1,x_2,\vdots,x_\Omega}
    \end{center}
    定义 $\mathrm{R}^{\Omega}$ 为\textbf{支付空间}（payoff space），即所有可能支付的集合；其中的向量
    \[
        x=\begin{bmatrix}x_1 \\ x_2 \\ \vdots \\ x_\Omega\end{bmatrix}
    \]
    就定义了一只证券。在定义证券时，我们未考虑任何在 0 期的可能支付；到期以前（在我们的框架中就是 0 期）证券产生的支付通常叫做\textbf{股利}（dividend）。
    
    在 Lucas 树例子中，把代表树的所有权的证券叫做\textbf{股票}（stock），它代表对实物资产（树）所产生的支付的要求权；这样的证券也叫做\textbf{金融资产}（financial asset）或\textbf{资产}，其支付来自\textbf{标的}（underlying）实物资产。

    \textbf{无风险债券}（riskless bond）是这样一种证券，它在所有状态下的支付都相等；若各状态下的支付均为 $1$，则其支付向量为 $\begin{bmatrix}1 \\ \vdots \\ 1\end{bmatrix}$。
    
    \textbf{风险证券}（如股票）的支付则随状态变化，例如支付向量为 $\begin{bmatrix}x_1 \\ \vdots \\ x_\Omega\end{bmatrix}$。
    
    金融证券具有三个基本特征：
    \begin{enumerate}
        \item 支付只依赖于未来的经济状况；
        \item 支付是外生给定的，不受参与者行为的影响；
        \item 状态集合及其概率分布是所有参与者的共同知识。
    \end{enumerate}

    \item \textbf{市场结构：}假设市场上有 $N$ 只可交易证券，用 $n=1,\dots,N$ 来标识。记第 $n$ 只证券的支付向量为
    \[
        X_n=\begin{bmatrix}x_{1,n} \\ x_{2,n} \\ \vdots \\ x_{\Omega,n}\end{bmatrix},
    \]
    把所有证券的支付向量并排，就得到 $\Omega\times N$ 的\textbf{市场结构矩阵}（也称支付矩阵）
    \[
        X\equiv[X_{\cdot,1},\,\cdots\,X_{\cdot,n},\,\cdots\,X_{\cdot,N}]=
        \begin{bmatrix}
            x_{1,1}   & x_{1,2}   & \cdots & x_{1,N}   \\
            \vdots    & \vdots    & \ddots & \vdots    \\
            x_{\omega,1}   & x_{\omega,2}   & \cdots & x_{\omega,N}   \\
            \vdots    & \vdots    & \ddots & \vdots    \\
            x_{\Omega,1} & x_{\Omega,2} & \cdots & x_{\Omega,N}
        \end{bmatrix}\equiv
        \begin{bmatrix}
            X_{1,\cdot} \\ \vdots \\ X_{\omega,\cdot} \\ \vdots \\ X_{\Omega,\cdot}
        \end{bmatrix}
    \]
    矩阵 $X$ 的每一列是一只证券在各状态下的支付（列向量），每一行是一个状态下各证券的支付（行向量）。通常有 $\Omega>N$。

    \item \textbf{证券组合：}对每只证券的持有数量构成\textbf{证券组合}（security portfolio），简称\textbf{组合}（portfolio）。若 $\theta_n$ 为证券 $n$ 的持有数量，则\textbf{组合}由向量$\theta=\begin{bmatrix}\theta_1 \\ \theta_2 \\ \vdots \\ \theta_N\end{bmatrix}$ 定义；该组合在 1 期的支付向量为 $X\theta$。对每个参与者 $k$，记其交易前的\textbf{初始证券组合}为$\bar{\theta}_k=\begin{bmatrix}\bar{\theta}_{k,1} \\ \bar{\theta}_{k,2} \\ \vdots \\ \bar{\theta}_{k,N}\end{bmatrix}$，所有参与者初始证券组合之和称为\textbf{市场组合}（market portfolio）$\theta_M$，即
    \[
        \theta_M=\sum_{k=1}^{K}\bar{\theta}_k,
    \]
    它表示市场上所有可交易证券的总供给。

    \item \textbf{市场化：}若支付 $x\in\mathrm{R}^{\Omega}$ 能由某交易组合复制出来，即存在 $\theta\in\mathrm{R}^{N}$ 使得 $X\theta=x$（有解），则称 $x$ 是\textbf{市场化的}（marketed）。所有市场化支付的集合为
    \[
        M=\{X\theta:\,\theta\in\mathrm{R}^{N}\}.
    \]
    一般地，$M$ 是支付空间 $\mathrm{R}^{\Omega}$ 的一个子空间，其维数等于 $\mathrm{rank}(X)\le\min\{N,\Omega\}$，只有证券支付列线性无关时才等于 $N$。允许任意实数组合保证了这个集合是线性子空间；若有卖空或持仓限制，受限可达集合未必是子空间。

\begin{tcolorbox}[example,title=PPT 补充：股票与债券能生成哪些支付]
Lucas 树股票支付 $[200,50]^{\mathsf T}$，单位无风险债券支付 $[1,1]^{\mathsf T}$，所以
\[
 X=\begin{bmatrix}200&1\\50&1\end{bmatrix}.
\]
给定目标支付 $y=[y_a,y_b]^{\mathsf T}$，解 $X\theta=y$ 得
\[
 \theta_{\rm stock}=\frac{y_a-y_b}{150},\qquad
 \theta_{\rm bond}=\frac{4y_b-y_a}{3}.
\]
例如支付 $[1,0]^{\mathsf T}$ 可由 $1/150$ 份股票与 $-1/3$ 份债券复制。股票与债券共同生成全部两状态支付；仅有股票时，只能生成其支付列的倍数。复制成本由价格 $S^{\mathsf T}\theta$ 决定，不能只看发生概率。
\end{tcolorbox}

    \item \textbf{交易过程：}交易可以通过\textbf{集中化交易}（centralized exchange）或\textbf{柜台交易}（over-the-counter，OTC）完成；具体的交易安排称为\textbf{市场微观结构}（market micro-structure）。记 $S=[S_1,\ldots,S_N]^{\mathsf T}$ 为证券价格，$h_k(S)$ 为参与者的\textbf{交易后总持仓}，$\bar\theta_k$ 为初始持仓。总持仓出清条件为
\[
 \sum_k h_k(S)=\sum_k\bar\theta_k=\theta_M.
\]
若改用净交易量 $\theta_k(S)=h_k(S)-\bar\theta_k$，则完全等价地写成 $\sum_k\theta_k(S)=0$。不能把净买入需求之和与总供给直接相比。PPT 的开盘报价例子说明的是价格下的\textbf{意愿需求}：若老王拟净卖出 $200$ 股、老张拟净买入 $300$ 股，尚有 $100$ 股超额净需求，须由其他参与者或调整后的意愿补足；实际成交的买卖量总是相等。此处只是出清机制的教学示意，不是对交易所撮合程序或某只股票当前报价的描述。

    交易成本、头寸限制（如卖空约束）、个体交易对价格的影响以及税收等，都会妨碍交易，统称为\textbf{市场摩擦}。\textbf{无摩擦市场}满足以下假设：
    \begin{enumerate}
        \item 参与者可无成本地进入或退出证券市场；
        \item 没有交易成本；
        \item 对参与者的证券持有没有限制（允许卖空）；
        \item 单个参与者的交易不影响市场价格（价格接受者）；
        \item 没有税收。
    \end{enumerate}
    税费、买卖价差与融资约束会使可执行交易的净成本偏离简单标价；此处用无摩擦假设建立基准，不把 PPT 中列举的某一时期税率当作模型常数。摩擦可能解释机构与交易安排的作用，但不保证存在任何人都能无风险获利的机会。
\end{enumerate}
\end{tcolorbox}

\begin{tcolorbox}[knowledge,title=2.4 基本经济模型]
现在我们已经对将要分析的经济的基本模型有了完整的描述。在作进一步的讨论和分析以前，先作一个总结。

\begin{tcolorbox}[knowledge]
\textbf{定义 2.5}\quad 一个\textbf{经济}定义如下：
\begin{enumerate}
    \item 有两个时期 $0$ 和 $1$；在 1 期有 $\Omega$ 个可能状态，对这些状态有概率测度 $P$。经济中只有一种不可储存的商品。
    \item 经济中有 $K$ 个参与者（$k=1,\dots,K$）：
    \begin{enumerate}[label=(\alph*)]
        \item 每一参与者对未来状态发生的可能性都有\textbf{相同的信息}，由 $P$ 描述；
        \item 每一参与者具有\textbf{禀赋} $e_k\in\mathrm{R}^{1+\Omega}_{+}$；
        \item 每一参与者有定义于 $C=\mathrm{R}^{1+\Omega}_{+}$ 上、且满足公理 1$\sim$公理 3 的偏好（不满足性、连续性、凸性）。
    \end{enumerate}
    \item 有一个市场结构为 $X$ 的\textbf{无摩擦证券市场}。
\end{enumerate}
\end{tcolorbox}

\begin{tcolorbox}[knowledge]
\textbf{定义 2.6}\quad 在上面定义的经济中，如果所有参与者的 1 期禀赋都可表示为其初始证券组合的支付，即对每个 $k$ 有 
\[
e_{k,1}=X\bar{\theta}_k
\]
则我们称之为\textbf{证券市场经济}（security-market economy）。
\end{tcolorbox}

这要求\textbf{每个人的实际禀赋可市场化}，不要求所有想象得到的支付都可交易，因此与第三章的“完全市场”不同。例如两状态只有支付 $[1,1]^{\mathsf T}$ 的债券，若人人禀赋也是该向量的倍数，它仍是证券市场经济，但并不完全。

PPT 以农业收入与劳动收入说明：如果某人的未来禀赋不在 $X$ 的列空间内，就不能完全表示为现有证券组合。仅凭资产未直接挂牌不能判断它不可复制，仍须检查支付。证券化、农业风险合约等可能扩大可达支付集合，但不会自动消除所有收入风险或保证帕累托效率。

为简化起见，在以后的讨论中有时会局限于证券市场经济的情形。
\end{tcolorbox}

\begin{tcolorbox}[knowledge,title=2.5 市场均衡]
在证券市场中，资源配置通过以下两方面的相互作用来实现：
\begin{enumerate}
    \item 给定证券及其未来支付与当前价格，每个参与者选择最优证券持有量；
    \item 所有参与者的总需求影响证券价格；当需求等于供给时达到\textbf{均衡}（equilibrium），此时各参与者都选择了最优证券需求且市场出清。
\end{enumerate}

\tcbline

\begin{enumerate}
    \item \textbf{参与者各自优化：}记 $\theta=\begin{bmatrix}\theta_1 \\ \vdots \\ \theta_N\end{bmatrix}$ 为\textbf{净购买}的证券组合（相对初始持有量的交易量），$\theta_n>0$ 表示\textbf{净买入}，$\theta_n<0$ 表示\textbf{净卖出}；最终总持仓 $h_{k,n}$ 的正负才分别表示多头和空头。该组合在 0 期的市场价值（成本）为 $S^{\mathsf T}\theta$。对禀赋为 $e_k$ 的参与者，净购买 $\theta_k$ 后消费计划 $c_k=\begin{bmatrix}c_{k,0} \\ c_{k,1}\end{bmatrix}$ 为
    \[
        c_{k,0}=\underbrace{e_{k,0}-S^{\mathsf T}\theta_k}_{\text{第零期禀赋减去购买证券的支出}},\qquad
        c_{k,1}=\underbrace{e_{k,1}+X\theta_k}_{\text{第 1 期禀赋加上证券在第一期的支付}}.
    \]
    其中 $S^{\mathsf T}\theta_k>0$ 表示净支出、减少当前消费，$S^{\mathsf T}\theta_k<0$ 表示净收入。对购买正价格债券这一特例，前者是储蓄、后者是借款；对一般含风险资产的组合，应称交易支出或收入，不能只凭符号笼统叫储蓄或借贷。所有可达消费计划的集合称为\textbf{预算集}
    \[
        B(e_k,\{X,S\})=\Bigl\{c\ge 0\;\Big|\;
        c_{k,0}=e_{k,0}-S^{\mathsf T}\theta_k,\;
        c_{k,1}=e_{k,1}+X\theta_k,\;
        \theta_k\in\mathrm{R}^{N}\Bigr\}.
    \]
    参与者选择 $\theta_k$ 以最大化效用：
    \[
    \begin{aligned}
        &\max_{\theta_k} U_k(c_k) \\
        &\text{s.t.}\quad
        c_{k,0}=e_{k,0}-S^{\mathsf T}\theta_k,\;
        c_{k,1}=e_{k,1}+X\theta_k,\;
        c_{k,0},\,c_{k,1}\ge 0
    \end{aligned}
    \]
    等价地，
    \[
        \max_{c_k\in B(e_k,\{X,S\})} U_k(c_k).
    \]
    该问题的解给出参与者 $k$ 对证券的\textbf{需求函数} $\theta_k(e_k,S)$，它依赖于价格向量 $S$ 与禀赋 $e_k$。

    \item \textbf{市场出清：}为避免对初始持仓重复计数，可区分两种一致写法。令 $\varepsilon_{k,1}$ 为不含证券支付的直接禀赋，$h_k=\bar\theta_k+\theta_k$ 为总持仓，则
    \[
      c_{k,1}=\varepsilon_{k,1}+Xh_k
      =\underbrace{(\varepsilon_{k,1}+X\bar\theta_k)}_{e_{k,1}}+X\theta_k.
    \]
    本章使用右边写法：$e_{k,1}$ 已包括初始证券支付，$\theta_k$ 是净交易，不能再把 $X\bar\theta_k$ 加一次。用总持仓时出清为 $\sum_kh_k=\sum_k\bar\theta_k$，用净交易时出清为 $\sum_k\theta_k=0$。这只是改写账目，并没有把实际证券供给变成零。

    如此一来，\textbf{证券市场出清}要求有人净买 $\theta_k>0$，就要有人净卖 $\theta_k<0$。对所有参与者加总：净买入之和必须为 0，否则多出来的份额没人卖、或少掉的份额没人接，即净需求为零：
    \[
        \underbrace{\sum_{k=1}^{K}\theta_k(e_k,S)=0}_{\text{证券市场出清条件}}.
    \]
    满足该条件的 $S$ 即为\textbf{均衡价格}。由 $c_{k,0}=e_{k,0}-S^{\mathsf T}\theta_k(e_k,S)$，对 $k$ 求和得
    \[
        \underbrace{\underbrace{\sum_{k=1}^{K}c_{k,0}}_{\text{0 期商品总需求}}
        =\sum_{k=1}^{K}e_{k,0}-S^{\mathsf T}\Bigl(\sum_{k=1}^{K}\theta_k(e_k,S)\Bigr)
        =\underbrace{\sum_{k=1}^{K}e_{k,0}}_{\text{0 期总供给}}}_{\text{第0期商品市场的出清}}
    \]
    左边是第0 期商品的总需求，右边是第0期的总供给，表示的是第0期商品市场的出清。同理，对每个状态 $\omega=1,\dots,\Omega$ 有
    \[
        \sum_{k=1}^{K}c_{k,\omega}=\sum_{k=1}^{K}e_{k,\omega},
    \]
    即\textbf{商品市场出清}。合起来得到\textbf{逐日期、逐状态的资源可行性}（PPT 在此也称作 Walras 法则的出清含义）：
    \[
        \sum_{k=1}^{K}c_k=\sum_{k=1}^{K}e_k.
    \]
    

    求均衡时，\textbf{首先}对任意价格 $S$ 解出各参与者的最优证券组合 $\theta_k(e_k,S)$，\textbf{第二}用市场出清条件解出均衡价格 $S$。形式上，
    \[
        S=S\bigl(P;\,\{U_k,\,e_k,\,k=1,\dots,K\};\,X\bigr),
    \]
    这是模型内均衡价格对基本要素的依赖关系，不自动保证均衡存在或唯一，也不是现实中所有波动的完备解释。在生产技术、信息和制度均固定的本模型中分析比较静态，不能把任何政策变化都直接等同于偏好 $U$ 的变化。即证券价格取决于经济面临的风险 $P$、参与者的偏好 $U_k$ 与禀赋 $e_k$ 以及市场结构 $X$ 等\textbf{基本面}（fundamentals）或\textbf{本源}（primitives）：经济面临的风险，经济中参与者的偏好和禀赋以及证券市场的结构。
\end{enumerate}
\end{tcolorbox}

\begin{tcolorbox}[knowledge,title=2.6 最优性,breakable=false]
通过证券市场进行的资源配置是否已是\textbf{最好的可能配置}？参与者各自选择最优消费计划，并不意味着社会资源配置一定最优；若市场不完备，结果可能是\textbf{次优的}。

\begin{tcolorbox}[knowledge]
\textbf{定义 2.7}\quad 配置 $\{c_k,\,\forall k\}$ \textbf{Pareto 占优于}配置 $\{c_k',\,\forall k\}$，如果
\[
    \forall k,\quad U_k(c_k)\ge U_k(c_k'),
\]
且至少对一个参与者严格不等式成立。

这就是说，对比两套分配结果，如果新配置帕累托占优于旧配置，说明：新配置对所有参与者的效用都至少不比原来差，且至少对一个参与者的效用严格提高，且这种“好”不以“损人”为前提，这是一种\textbf{更具效率}的体现。
\end{tcolorbox}



\begin{tcolorbox}[knowledge]
\textbf{定义 2.8}\quad 给定经济的总供给 $\{e_k,\,\forall k\}$，配置 $\{c_k,\,\forall k\}$ 是\textbf{可行的}（feasible），如果
\[
    \underbrace{\sum_{k=1}^{K}c_k=\sum_{k=1}^{K}e_k}_{\text{说明所有状态下所有人的消费之和等于总禀赋之和，没有闲置}}.
\]
\end{tcolorbox}

\begin{tcolorbox}[knowledge]
\textbf{定义 2.9}\quad 若配置 $\{c_k,\,\forall k\}$ 可行，且不存在其他可行配置 Pareto 占优于它，则该配置是\textbf{Pareto 最优的}（Pareto optimal）。Pareto 最优即在不牺牲其他参与者福利、也不使用更多资源的前提下，无法改进任一参与者的福利（某人福利的提高，必定以损害他人的利益为前提）。给定偏好后，通常只有可行配置的一个子集是最优的。
\end{tcolorbox}

这里比较的是人的\textbf{整段消费计划的效用}，并不要求每个日期、每个状态下的消费都逐项增加。

Pareto 最优配置也叫做\textbf{有效}（efficient）配置。若证券市场使参与者达到 Pareto 最优配置，则称该市场在\textbf{资源配置意义上有效}。这与“价格反映信息”的有效市场假说是不同概念；效率也不等于平均分配或公平。
\end{tcolorbox}


\begin{tcolorbox}[example,title=例题：帕累托最优]
    \begin{question}
        \textbf{例 2.5}\quad 考虑一个两期、1 期有两种状态 $a$、$b$ 的经济。经济总禀赋：
        \begin{center}
        \WJEventTree[left={$100$},stem]{400,25}
        \end{center}
        有两个参与者，效用函数均为 $u(c_0,c_{1a},c_{1b})=\sqrt{c_0}+\tfrac12(\sqrt{c_{1a}}+\sqrt{c_{1b}})$（$c=\begin{bmatrix}c_0 \\ c_{1a} \\ c_{1b}\end{bmatrix}$）。树中每个节点括号内为 $(c_{1,\cdot},\,c_{2,\cdot})$：该节点上参与者 1、2 的消费量（不是整段消费计划列向量）。考虑三种配置：
        \begin{center}
        \begin{tabular}{@{}c@{\quad}c@{\quad}c@{}}
            \textbf{配置 A} & \textbf{配置 B} & \textbf{配置 C} \\[0.35em]
            \WJEventTree[left={$(100,0)$},stem]{ {(0,400)}, {(0,25)} } &
            \WJEventTree[left={$(50,50)$},stem]{ {(200,200)}, {(\tfrac{25}{2},\tfrac{25}{2})} } &
            \WJEventTree[left={$(100,0)$},stem]{ {(400,0)}, {(25,0)} }
        \end{tabular}
        \end{center}
        配置 A 中参与者 1 在 0 期消费全部、参与者 2 在 1 期消费全部；B 为两参与者在各期、各状态均分；C 中参与者 1 始终消费全部。哪种配置是 Pareto 最优的？（完整消费计划 $c_1$、$c_2$ 见解答。）
    \end{question}
    \begin{solution}
        与前文及习题 2.2 相同，参与者 $i$ 的\textbf{消费计划}是 $\mathrm{R}^{3}_{+}$ 中的列向量（两状态记为 $a,b$）：
        \[
            c_i=\begin{bmatrix}c_{i,0} \\ c_{i,1a} \\ c_{i,1b}\end{bmatrix},
            \qquad U_i(c_i)=\sqrt{c_{i,0}}+\tfrac12\bigl(\sqrt{c_{i,1a}}+\sqrt{c_{i,1b}}\bigr).
        \]
        题中事件树每个节点上的 $(\cdot,\cdot)$ 只是该节点上\textbf{参与者 1、2 的消费量}；整段计划仍用上面的 $c_1$、$c_2$ 表示。

        \textbf{可行性:}三个配置合起来在每一日期—状态下满足 $c_{1,\cdot}+c_{2,\cdot}$ 等于总禀赋（0 期 $100$，状态 $a$ 为 $400$，状态 $b$ 为 $25$），故 A、B、C 均为可行配置。

        \textbf{各配置的效用：}
        \begin{itemize}
            \item \textbf{配置 A：}
            \[
                c_1=\begin{bmatrix}100 \\ 0 \\ 0\end{bmatrix},\qquad
                c_2=\begin{bmatrix}0 \\ 400 \\ 25\end{bmatrix}.
            \]
            \[
                \begin{aligned}
                    U_1(c_1) &= \sqrt{100}+\tfrac12(0+0)=10,\qquad 
                    U_2(c_2) &= 0+\tfrac12(\sqrt{400}+\sqrt{25})
                    =\tfrac12(20+5)=12.5.
                \end{aligned}
            \]
            效用对为 $(U_1,U_2)=(10,\,12.5)$。

            \item \textbf{配置 B}（均分）：$c_1=c_2=
            \begin{bmatrix}50 \\ 200 \\ \tfrac{25}{2}\end{bmatrix}$。
            \[
                U_1(c_1)=U_2(c_2)=\sqrt{50}+\tfrac12\Bigl(\sqrt{200}+\sqrt{\tfrac{25}{2}}\Bigr)
                \approx 7.071+8.839=15.91.
            \]
            效用对为 $(15.91,\,15.91)$。

            \item \textbf{配置 C：}
            \[
                c_1=\begin{bmatrix}100 \\ 400 \\ 25\end{bmatrix},\qquad
                c_2=\begin{bmatrix}0 \\ 0 \\ 0\end{bmatrix}.
            \]
            \[
                \begin{aligned}
                    U_1(c_1) = \sqrt{100}+\tfrac12(\sqrt{400}+\sqrt{25})=10+12.5=22.5, \quad
                    U_2(c_2) = 0.
                \end{aligned}
            \]
            效用对为 $(22.5,\,0)$。
        \end{itemize}

        \textbf{Pareto 占优与 Pareto 最优。}
        \begin{itemize}
            \item 比较 A 与 B：$U_1(B)=15.91>10=U_1(A)$ 且 $U_2(B)=15.91>12.5=U_2(A)$，两人都严格更好，故\textbf{配置 B Pareto 占优于配置 A}。因此 A \textbf{不是} Pareto 最优。

            \item 比较 B 与 C：$U_1(C)=22.5>U_1(B)$，但 $U_2(C)=0<U_2(B)$，有人变差，故 C 不 Pareto 占优 B，B 也不 Pareto 占优 C。

            \item \textbf{配置 B 是 Pareto 最优的}：仅比较 A、B、C 还不够，须排除\textbf{所有}可行改善。对任一节点总资源 $E$ 和分配 $x,E-x$，凹性给出 $\sqrt x+\sqrt{E-x}\le2\sqrt{E/2}$。将 0 期与两状态分别加权相加，得任何可行配置的 $U_1+U_2$ 都不超过均分 B 的总效用。若有人严格改善、另一人不变差，总效用必上升，矛盾。因此 B 最优；这里加总效用只是在给定这两个效用函数后证明结论的数学工具，不是一般的人际福利比较标准。

            \item \textbf{配置 C 是 Pareto 最优的}：若要提高参与者 2 的效用（必须给 2 正的消费），只能从参与者 1 已占用的 $100$、$400$、$25$ 中让出，从而必使 $U_1$ 下降，无法在不损害 1 的情况下改进 2。
        \end{itemize}

        
    \end{solution}

    \begin{tcolorbox}[warning,title=注意]
        \textbf{结论：}A 非 Pareto 最优；\textbf{B 与 C 都是} Pareto 最优配置（同一总禀赋下不同的有效分配）。可以看出：
        \begin{itemize}
            \item 在同一可行集（同一总禀赋、满足 $\sum_k c_k=\sum_k e_k$）里，Pareto 最优配置往往不止一个。
            \item 更能看到，帕累托最优仅考虑了“效率”，但没考虑“公平”，例如配置 C 中参与者 2 的效用为0，但仍属于帕累托最优。
        \end{itemize}
        \end{tcolorbox}
\end{tcolorbox}


\begin{tcolorbox}[knowledge,title=2.7 本章小结]
本章完成了在市场经济中通过资源配置分析所使用的基本框架的构建。资源配置是通过证券市场完成的：交易证券的价格引导参与者通过交易达到他们希望的配置，而交易又最终决定了证券的市场价格。证券价格在引导个体参与者和平衡整个经济方面都起着关键作用。因此，我们关注证券价格如何决定，这也称为\textbf{资产定价}（asset pricing）。

\begin{enumerate}
    \item 从金融的角度看，一个经济的结构包括三个方面：（1）它所处的\textbf{自然环境}；（2）\textbf{经济的参与者}；（3）\textbf{金融市场}。

    \item \textbf{外部环境}包括两个要素：\textbf{时间}和\textbf{风险}。为刻画时间，假设有 0 和 1 两个时期；为刻画风险，假设 1 期有 $\Omega$ 个可能状态，用概率测度 $P$ 描述。经济中只有一种不可储存的商品。

    \item 经济中有 $K$ 个\textbf{参与者}。每个参与者的经济特征包括其经济资源和经济需求。具体假设为：
    \begin{enumerate}[label=(\alph*)]
        \item 每个参与者 $k$ 具有禀赋 $e_k\in\mathrm{R}^{1+\Omega}_{+}$；
        \item 每个参与者对未来状态发生的可能性拥有相同的信息，完全由概率测度 $P$ 描述；
        \item 尽管参与者给定其禀赋和信息后可能拥有生产技术，这里先假设他们不拥有任何生产技术；
        \item 每个参与者在 $C=\mathrm{R}^{1+\Omega}_{+}$ 上具有偏好，且满足不满足性、连续性和凸性公理；
        \item 给定这些偏好关系，可用效用函数来描述。
    \end{enumerate}

    \item \textbf{金融市场}由一组证券构成：
    \begin{enumerate}[label=(\alph*)]
        \item 它们的支付只依赖于未来实际的经济状态；
        \item 它们的支付是外生给定的，且不受经济参与者行为的影响；
        \item 参与者知道且一致认同证券支付的概率分布；
        \item 市场结构为 $X$，且市场中无摩擦。
    \end{enumerate}

    \item \textbf{市场均衡}必须满足两个条件：
    \begin{enumerate}[label=(\alph*)]
        \item 每个参与者各自最优化其证券需求量；
        \item \textbf{市场出清}，即交易证券的总需求等于总供给，此时商品市场也是出清的。
    \end{enumerate}

    \item 如果配置 $\{c_k,\,\forall k\}$ 是可行的，且不存在另外的可行配置 Pareto 占优于它，则该配置是 \textbf{Pareto 最优的}。如果证券市场允许参与者达到 Pareto 最优配置，就称它为\textbf{有效市场}。

    \item \textbf{对数效用下的支出份额。}设 \(U=\log c_0+\rho\log c_1\)，两期且 1 期只有一个状态，只有收益率为 \(r_F\) 的无风险证券。以 0 期消费计价的财富为 \(w=e_0+e_1/(1+r_F)\)。最优消费满足
    \[
        c_1^{*}=\rho(1+r_F)\,c_0^{*},
    \]
    因而 0 期消费占财富的份额只由 \(\rho\) 决定：
    \[
        c_0^{*}=\frac{w}{1+\rho},\qquad
        c_1^{*}=\frac{\rho(1+r_F)}{1+\rho}\,w.
    \]
    \(c_0^{*}/w=1/(1+\rho)\) 不随 \(r_F\) 改变；利率只通过 \(w\) 里未来禀赋的现值影响 \(c_0^{*}\)。这是对数效用（相对风险厌恶为 1）的结论。对幂效用 $u(c)=c^{1-\gamma}/(1-\gamma)$，一阶条件为 $c_1/c_0=[\rho(1+r_F)]^{1/\gamma}$，所以
    \[
      \frac{c_0^*}{w}=\frac1{1+\rho^{1/\gamma}(1+r_F)^{1/\gamma-1}}.
    \]
    幂效用的份额仍不依赖财富水平，但一般会随利率改变；对数情形 $\gamma=1$ 才消去上式的利率项。负指数效用一般也不具有同样的固定份额。
\end{enumerate}
\end{tcolorbox}
\fi

% !TeX root = ../../金融经济学_王江_讲义及习题.tex
% 02_基本框架 课后习题（由章节文件拆分）

\begin{problem}
\problemno{2.1} $U(c)$ 和 $V(c)$ 是两个效用函数，$c\in\mathrm{R}_{+}^{n}$，且 $V(x)=f(U(x))$，其中 $f(\cdot)$ 是一正单调函数。证明这两个效用函数表示了相同的偏好。
\end{problem}
\begin{solution}
定义 2.4：效用函数表示偏好 $\succeq$，是指对任意两个消费计划，$U(c_{1})\ge U(c_{2})$ 当且仅当 $c_{1}\succeq c_{2}$。两个效用函数表示相同偏好，就是它们给出完全一样的排序。正单调变换不改变不等号的方向，所以 $V=f(U)$ 与 $U$ 排序相同。

假设 $U(c)$ 表示的偏好关系为 $\succeq$，那么 $\forall\, c_{1},\,c_{2}\in\mathrm{R}_{+}^{N}$ 有
\[
    U(c_{1})\ge U(c_{2})\Leftrightarrow c_{1}\succeq c_{2}
\]
而 $f(\cdot)$ 是正单调函数，因而
\[
    V(c_{1})=f(U(c_{1}))\ge f(U(c_{2}))=V(c_{2})\Leftrightarrow U(c_{1})\ge U(c_{2})
\]
因此 $V(c_{1})\ge V(c_{2})\Leftrightarrow c_{1}\succeq c_{2}$，即 $V(c)$ 表示的偏好也是 $\succeq$。
\end{solution}

\begin{problem}
\problemno{2.2*} 在 1 期，经济有两个可能状态 $a$ 和 $b$，它们的发生概率相等：
\begin{center}
\WJEventTree{a,b}
\end{center}
考虑定义在消费计划 $c=\begin{bmatrix}c_{0} \\ c_{1a} \\ c_{1b}\end{bmatrix}$ 上的效用函数：
\[
    \begin{aligned}
        U(c) &= \log c_{0}+\tfrac12(\log c_{1a}+\log c_{1b})\\
        U(c) &= \tfrac{1}{1-\gamma}c_{0}^{1-\gamma}+\tfrac12\Bigl(\tfrac{1}{1-\gamma}c_{1a}^{1-\gamma}+\tfrac{1}{1-\gamma}c_{1b}^{1-\gamma}\Bigr)\\
        U(c) &= -e^{-ac_{0}}-\tfrac12\bigl(e^{-ac_{1a}}+e^{-ac_{1b}}\bigr)
    \end{aligned}
\]
证明它们满足：不满足性、连续性和凸性。
\end{problem}
\begin{solution}
三个效用函数都具有形式 $U(c)=u(c_{0})+\tfrac12\bigl(u(c_{1a})+u(c_{1b})\bigr)$。对数与幂效用要求底数为正，指数效用 $u(x)=-e^{-ax}$ 在 $a>0$ 时于整个实数轴上有定义。下面统一在严格正消费集
\[
    C=\bigl\{c\in\mathrm{R}^{3}:c_{0}>0,\;c_{1a}>0,\;c_{1b}>0\bigr\}
\]

上证明；并设幂效用中 $\gamma>0$、$\gamma\neq 1$，指数效用中 $a>0$。对应的期效用为
\[
    u(x)=\log x,\qquad
    u(x)=\frac{x^{1-\gamma}}{1-\gamma},\qquad
    u(x)=-e^{-ax},\qquad x>0.
\]

三者都满足 $u'(x)>0$、$u''(x)<0$，因而严格递增且严格凹：
\[
    (\log x)'=\frac1x,\quad (\log x)''=-\frac1{x^{2}},
\]
\[
    \Bigl(\frac{x^{1-\gamma}}{1-\gamma}\Bigr)'=x^{-\gamma},\quad
    \Bigl(\frac{x^{1-\gamma}}{1-\gamma}\Bigr)''=-\gamma x^{-\gamma-1},
\]
\[
    (-e^{-ax})'=ae^{-ax},\quad (-e^{-ax})''=-a^{2}e^{-ax}.
\]

\begin{enumerate}[label=(\alph*)]
    \item 不满足性（喜多厌少）。设 $c,c'\in C$ 且 $c\ge c'$。由 $u$ 在 $(0,\infty)$ 上递增，$u(c_{i})\ge u(c_{i}')$，故 $U(c)\ge U(c')$。
    \item 连续性。要证：若 $\{c^{(n)}\}\subset C$ 且 \(\lim_{n\to\infty}c^{(n)}=c\in C\)，则
    \[
        \lim_{n\to\infty}U\bigl(c^{(n)}\bigr)=U(c).
    \]

    按分量有 \(\lim_{n\to\infty}c_{i}^{(n)}=c_{i}>0\)。\(u\) 在 \((0,\infty)\) 上连续，因而在点 \(c_{i}\) 连续。由连续的序列定义，极限可以代入 \(u\)：
    \[
        \lim_{n\to\infty}u\bigl(c_{i}^{(n)}\bigr)=u\Bigl(\lim_{n\to\infty}c_{i}^{(n)}\Bigr)=u(c_{i}).
    \]
    $U$ 只是三项 $u$ 的和，极限穿过和式：
    \[
        \begin{aligned}
            \lim_{n\to\infty}U\bigl(c^{(n)}\bigr)
            &=\lim_{n\to\infty}u\bigl(c_{0}^{(n)}\bigr)
            +\tfrac12\lim_{n\to\infty}u\bigl(c_{1a}^{(n)}\bigr)
            +\tfrac12\lim_{n\to\infty}u\bigl(c_{1b}^{(n)}\bigr)\\
            &=u(c_{0})+\tfrac12\bigl(u(c_{1a})+u(c_{1b})\bigr)
            =U(c).
        \end{aligned}
    \]
    \item 凸性。$u$ 严格凹，故对 $\alpha\in(0,1)$ 与 $c,c'\in C$，
    \[
        u\bigl(\alpha c_{i}+(1-\alpha)c_{i}'\bigr)\ge\alpha u(c_{i})+(1-\alpha)u(c_{i}'),
    \]
    且凸组合仍在 $C$ 中。因此
    \[
        U\bigl(\alpha c+(1-\alpha)c'\bigr)\ge\alpha U(c)+(1-\alpha)U(c').
    \]

    若 $U(c)>U(c')$，则右端严格大于 $U(c')$，即 $U\bigl(\alpha c+(1-\alpha)c'\bigr)>U(c')$。
\end{enumerate}
\end{solution}

\begin{problem}
\problemno{2.3} $U(c)=c-\tfrac12 a c^{2}$ 是一可能的效用函数，其中 $c\in\mathrm{R}_{+}$，$a$ 是非负的系数。$U(c)$ 具有不满足性吗？如果不是，那么 $a$ 取什么值和/或 $c$ 在什么范围内时 $U(c)$ 具有不满足性？
\end{problem}
\begin{solution}
不一定，比如当 $a=1$ 时，$U(\tfrac12)=\tfrac38<U(1)=\tfrac12>U(3)=-1.5$。$U(c)$ 不具有不满足性。

当 $a=0$ 时，$U(c)=c$ 具有不满足性；

当 $a>0$ 时，当 $c\in[0,\tfrac{1}{a}]$ 时 $U(c)$ 具有不满足性。

\begin{center}
\footnotesize
\begin{tabular}{@{}ccc@{}}
\begin{tikzpicture}[baseline=(current bounding box.center), x=1cm, y=1cm]
    \path (-0.32,-1.55) rectangle (2.55,1.62);
    \draw[-{Stealth[length=1.5mm]}, thin] (0,0) -- (2.28,0) node[below left, inner sep=1pt] {$c$};
    \draw[-{Stealth[length=1.5mm]}, thin] (0,0) -- (0,1.48) node[above left, inner sep=1pt] {$U$};
    \draw[blue!65!black, thick, domain=0:2] plot (\x,{0.65*\x});
\end{tikzpicture}
&
\begin{tikzpicture}[baseline=(current bounding box.center), x=1cm, y=1cm]
    \path (-0.32,-1.55) rectangle (2.55,1.62);
    \draw[-{Stealth[length=1.5mm]}, thin] (0,0) -- (2.28,0) node[below left, inner sep=1pt] {$c$};
    \draw[-{Stealth[length=1.5mm]}, thin] (0,0) -- (0,1.48) node[above left, inner sep=1pt] {$U$};
    \draw[blue!65!black, thick, domain=0:1, samples=40, smooth] plot ({2*\x},{2.2*(\x-0.5*\x*\x)});
    \draw[dashed, black!45] (2,0) -- (2,1.1);
    \node[below, inner sep=1pt] at (2,0) {$1/a$};
\end{tikzpicture}
&
\begin{tikzpicture}[baseline=(current bounding box.center), x=1cm, y=1cm]
    \path (-0.32,-1.55) rectangle (2.55,1.62);
    \draw[-{Stealth[length=1.5mm]}, thin] (0,0) -- (2.35,0) node[above right, inner sep=1pt] {$c$};
    \draw[-{Stealth[length=1.5mm]}, thin] (0,-1.4) -- (0,0.85) node[above left, inner sep=1pt] {$U$};
    \draw[blue!65!black, thick, domain=0:3.05, samples=50, smooth] plot ({0.68*\x},{0.85*(\x-0.5*\x*\x)});
    \draw[dashed, black!45] (0.34,0) -- (0.34,0.319);
    \draw[dashed, black!45] (0.68,0) -- (0.68,0.425);
    \draw[dashed, black!45] (2.04,0) -- (2.04,-1.275);
    \fill (0.34,0.319) circle (1pt);
    \fill (0.68,0.425) circle (1pt);
    \fill (2.04,-1.275) circle (1pt);
    \draw (0.34,0.045) -- (0.34,-0.045) node[below, inner sep=1.5pt] {$\tfrac12$};
    \draw (0.68,0.045) -- (0.68,-0.045) node[below, inner sep=1.5pt] {$1$};
    \draw (2.04,0.045) -- (2.04,-0.045) node[below, inner sep=1.5pt] {$3$};
\end{tikzpicture}
\\[0.15em]
$a=0$，$U(c)=c$ & $a>0$，$c\in[0,1/a]$ & $a=1$，$U(\tfrac12)<U(1)>U(3)$
\end{tabular}
\end{center}
\end{solution}

\begin{problem}
\problemno{2.4} 考虑一个经济，它在 1 期有三个可能状态：$a$，$b$ 和 $c$：
\begin{center}
\WJEventTree{a,b,c}
\end{center}
证券市场包括证券 1 和 2，它们具有如下的支付向量：$X_{1}=\begin{bmatrix}1 \\ 1 \\ 1\end{bmatrix}$ 以及 $X_{2}=\begin{bmatrix}1 \\ 2 \\ 3\end{bmatrix}$。它们的价格分别为 $S_{1}$ 和 $S_{2}$。
\begin{enumerate}[label=(\alph*)]
    \item 描述这个经济的支付空间。
    \item 写出这个经济的市场结构矩阵 $X$。
    \item 考虑含有 $\theta_{1}$ 单位的证券 1 和 $\theta_{2}$ 单位的证券 2 的组合。写出这个组合的支付向量。这个组合的价格是多少？
    \item 假设这个市场中总共有 $K$ 个参与者。每个参与者的禀赋是 1 单位的证券 1 和 2 单位的证券 2。这时的市场组合是什么？市场组合的支付向量是什么？市场组合的总价值是多少？
    \item 写出市场化支付的集合。
    \item 如果市场不允许卖空。市场化支付的集合是什么？
    \item 现在引入新的证券 3，它的支付向量为 $X_{3}=\begin{bmatrix}0 \\ 0 \\ 1\end{bmatrix}$。写出新的市场结构矩阵。在这个市场结构下，市场化支付集合是什么？
\end{enumerate}
\end{problem}
\begin{solution}
三个状态 $a,b,c$ 对应 $\Omega=3$。记
\[
    X_1=\begin{bmatrix}1 \\ 1 \\ 1\end{bmatrix},\quad
    X_2=\begin{bmatrix}1 \\ 2 \\ 3\end{bmatrix},\quad
    \theta=\begin{bmatrix}\theta_1 \\ \theta_2\end{bmatrix}.
\]

组合在 1 期的支付向量为 $X_1\theta_1+X_2\theta_2$，也可写成矩阵乘积 $X\theta$，其中 $X=[X_1,\,X_2]$。

\begin{enumerate}[label=(\alph*)]
    \item \textbf{支付空间}：所有可能在三个状态下实现的支付向量的集合，即 $\mathrm{R}^{\Omega}=\mathrm{R}^{3}$（每个分量对应状态 $a,b,c$ 下的支付）。

    \item \textbf{市场结构矩阵}：两列分别为两只证券的支付向量，
    \[
        X=\begin{bmatrix}1&1\\1&2\\1&3\end{bmatrix}=[X_1,\,X_2].
    \]

    \item \textbf{组合支付与价格：}持有 $(\theta_1,\theta_2)$ 时，
    \[
        X\theta=\theta_1 X_1+\theta_2 X_2=
        \begin{bmatrix}\theta_1+\theta_2 \\ \theta_1+2\theta_2 \\ \theta_1+3\theta_2\end{bmatrix}.
    \]

    0 期购买成本（组合价格）为 $\theta_1 S_1+\theta_2 S_2$。

    \item \textbf{市场组合：}$K$ 个参与者每人禀赋为 $1$ 单位证券 1、$2$ 单位证券 2，加总得市场组合
    \[
        \theta_M=\begin{bmatrix}K \\ 2K\end{bmatrix}.
    \]

    其支付向量为
    \[
        X\theta_M=K X_1+2K X_2=
        \begin{bmatrix}3K \\ 5K \\ 7K\end{bmatrix}.
    \]

    总价值为 $K S_1+2K S_2$。

    \item \textbf{市场化支付（允许卖空）：}能由证券组合复制的 1 期支付集合为
    \[
        M=\bigl\{Y\in\mathrm{R}^{3}: Y=X\theta=[X_1,\,X_2]\begin{bmatrix}\theta_1 \\ \theta_2\end{bmatrix}=\theta_1 X_1+\theta_2 X_2,\;\theta_1,\theta_2\in\mathrm{R}\bigr\},
    \]

    即 $X_1,X_2$ 张成的二维子空间（只有两只证券，维度 $\le 2$）。

    \item \textbf{禁止卖空：}仅允许 $\theta_1,\theta_2\ge 0$，为 $X_1,X_2$ 的\textbf{非负}线性组合：
    \[
        M_{+}=\bigl\{Y\in\mathrm{R}^{3}: Y=X\theta=[X_1,\,X_2]\begin{bmatrix}\theta_1 \\ \theta_2\end{bmatrix}=\theta_1 X_1+\theta_2 X_2,\;\theta_1,\theta_2\in\mathrm{R}_{+}\bigr\}.
    \]

    \item \textbf{加入证券 3：}其支付向量为 $X_3=\begin{bmatrix}0 \\ 0 \\ 1\end{bmatrix}$，仅在状态 $c$ 支付 $1$。两只证券时的市场结构是 $X=[X_1,\,X_2]$。加入证券 3 后的市场结构另记为 $X'$，以免与 $X$ 相混：
    \[
        X'=\begin{bmatrix}1&1&0\\1&2&0\\1&3&1\end{bmatrix}=[X_1,\,X_2,\,X_3].
    \]

    三列线性无关，可表示任意 $Y\in\mathrm{R}^{3}$，故（允许卖空时）
    \[
        M=\bigl\{Y\in\mathrm{R}^{3}: Y=\theta_1 X_1+\theta_2 X_2+\theta_3 X_3,\;\theta_1,\theta_2,\theta_3\in\mathrm{R}\bigr\}=\mathrm{R}^{3}.
    \]
\end{enumerate}
\end{solution}

\begin{problem}
\problemno{2.5} 在练习 2.4 中定义的只存在证券 1 和 2 的经济中。考虑一个禀赋为 $\theta_{1}$ 单位的证券 1 和 $\theta_{2}$ 单位的证券 2 的参与者。写出他的预算集。
\end{problem}
\begin{solution}
无摩擦，允许卖空。预算集是花费不超过禀赋市值的那些 1 期支付：
\[
    B_{12}=\Bigl\{C\in\mathrm{R}_{+}^{3}:
    C=\alpha_{1}X_{1}+\alpha_{2}X_{2},\;
    \alpha_{1}S_{1}+\alpha_{2}S_{2}\le\theta_{1}S_{1}+\theta_{2}S_{2},\;
    \alpha_{1},\alpha_{2}\in\mathrm{R}\Bigr\}.
\]
\end{solution}

\begin{problem}
\problemno{2.6} 在上面的练习中引入练习 2.4 中定义的证券 3，它的价格为 $S_{3}$。这时，参与者的预算集是什么（他在证券 3 上的禀赋为 0）？证明由证券 1、2、3 构成的预算集包含仅由证券 1、2 构成的预算集。
\end{problem}
\begin{solution}
证券 3 上的禀赋为 $0$，可交易财富仍是 $\theta_{1}S_{1}+\theta_{2}S_{2}$。预算集为
\[
    B_{123}=\Bigl\{C\in\mathrm{R}_{+}^{3}:
    C=\alpha_{1}X_{1}+\alpha_{2}X_{2}+\alpha_{3}X_{3},\;
    \alpha_{1}S_{1}+\alpha_{2}S_{2}+\alpha_{3}S_{3}\le\theta_{1}S_{1}+\theta_{2}S_{2},\;
    \alpha_{1},\alpha_{2},\alpha_{3}\in\mathrm{R}\Bigr\}.
\]

练习 2.5 的预算集是
\[
    B_{12}=\Bigl\{C\in\mathrm{R}_{+}^{3}:
    C=\alpha_{1}X_{1}+\alpha_{2}X_{2},\;
    \alpha_{1}S_{1}+\alpha_{2}S_{2}\le\theta_{1}S_{1}+\theta_{2}S_{2},\;
    \alpha_{1},\alpha_{2}\in\mathrm{R}\Bigr\}.
\]

设 $C\in B_{12}$，对应头寸为 $\alpha_{1},\alpha_{2}$。取 $\alpha_{3}=0$，则
\[
    C=\alpha_{1}X_{1}+\alpha_{2}X_{2}+\alpha_{3}X_{3},
\]

并且
\[
    \alpha_{1}S_{1}+\alpha_{2}S_{2}+\alpha_{3}S_{3}=\alpha_{1}S_{1}+\alpha_{2}S_{2}\le\theta_{1}S_{1}+\theta_{2}S_{2}.
\]

故 $C\in B_{123}$，即 $B_{12}\subseteq B_{123}$。
\end{solution}

\begin{problem}
\problemno{2.7*} 考虑一个在 1 期只有一个可能状态的经济。（在这种情况下不存在不确定性。）参与者 1 的 0 期禀赋为 100 而 1 期禀赋为 1，即他的禀赋向量为 $\begin{bmatrix}100 \\ 1\end{bmatrix}$。他的偏好可以表示成如下形式：
\[
    U(c_{0},c_{1})=\log c_{0}+\rho\log c_{1}.
\]

系数 $\rho$ 为反映参与者在当前消费和未来消费之间相对偏好的参数。有一只证券，它的 0 期价格为 1、1 期支付为 $1+r_{F}$。这里，$r_{F}$ 是利率。
\begin{enumerate}[label=(\alph*)]
    \item 如果这个参与者不能在市场上进行交易，那么他的消费计划以及相应的效用 $U^{a}$ 是什么？
    \item 现在假设他可以在市场上进行交易。
    \begin{itemize}
        \item 他的预算集是什么？以当前消费为单位，他的总财富 $w$ 是多少？
        \item 写出参与者的优化问题。令 $c_{0}$ 为参与者的当前（即 0 期）最优消费，$s$ 为最优储蓄以及 $U^{b}$ 为在最优策略下得到的效用。求解他的最优消费/储蓄选择以及相应的效用。把 $U^{b}$ 表示成财富 $w$、利率 $r_{F}$ 和偏好系数 $\rho$ 的函数。
        \item 讨论参与者的最优选择如何依赖于利率 $r_{F}$ 和偏好系数 $\rho$。给出解释。
    \end{itemize}
    \item 证明 $U^{b}\ge U^{a}$。
    \item 令 $g$ 为参与者由于能够在证券市场上交易而获得的益处。它的定义为
    \[
        U^{b}(w-g)=U^{a}.
    \]

    计算 $g$。讨论 $g$ 如何依赖于 $\rho$？$g$ 如何依赖于 $r_{F}$？给出解释。
\end{enumerate}
\end{problem}
\begin{solution}
禀赋 $e_0=100$、$e_1=1$。无风险证券：0 期价格 $1$，1 期支付 $1+r_F$。记 $S$ 为 0 期\textbf{储蓄}（买入该证券的金额；$S<0$ 表示借贷）。

\begin{enumerate}[label=(\alph*)]
    \item \textbf{不能交易（自给自足）：}只能消费禀赋，
    \[
        c_0=100,\qquad c_1=1,\qquad
        U^{a}=\log(100)+\rho\log(1)=\log 100.
    \]

    \item \textbf{可以交易。}
    \begin{itemize}
        \item \textbf{预算集与财富 $w$。}买卖证券等价于在两期之间转移消费：0 期少消费 $S$ 用于储蓄，1 期增加 $S(1+r_F)$。故
        \[
            c_0=100-S,\qquad c_1=1+S(1+r_F),\qquad S\in\mathrm{R}.
        \]

        以 0 期消费为计价单位，把 1 期禀赋折现到 0 期（用市场利率 $r_F$）：
        \[
            w=100+\frac{1}{1+r_F}.
        \]

        这是参与者按市价出售全部禀赋后可在 0 期支配的\textbf{总财富}（现值）。任意满足上述两式的 $(c_0,c_1)$ 都在预算集内。

        \item \textbf{优化问题。}在预算集上最大化对数效用：
        \[
            \max_{S\in\mathrm{R}}\ \log(100-S)+\rho\log\bigl[1+S(1+r_F)\bigr].
        \]

        一阶条件（边际效用替代率 = 跨期相对价格）：
        \[
            \underbrace{\frac{1}{100-S}}_{\text{0 期边际效用}}
            =\rho\underbrace{\frac{1+r_F}{1+S(1+r_F)}}_{\text{1 期边际效用}\times\text{价格比}}.
        \]

        整理得 $\bigl[1+S(1+r_F)\bigr]=\rho(1+r_F)(100-S)$，即
        \[
            1+S(1+r_F)(1+\rho)=100\rho(1+r_F).
        \]
        
        其解记为最优储蓄
        \[
            S^{*}=\frac{100\rho(1+r_F)-1}{(1+\rho)(1+r_F)}.
        \]

        预算约束是前面的两式 \(c_0=100-S\)、\(c_1=1+S(1+r_F)\)，其中 \(c_0,c_1,S\) 仍是可行选择，不是最优值。把 \(S^{*}\) 代回去，得到最优消费 \(c_0^{*},c_1^{*}\)：
        \[
            \begin{aligned}
                c_0^{*}&=100-S^{*}
                =100-\frac{100\rho(1+r_F)-1}{(1+\rho)(1+r_F)}
                =\frac{100(1+r_F)+1}{(1+\rho)(1+r_F)},\\[0.4em]
                c_1^{*}&=1+S^{*}(1+r_F)
                =1+(1+r_F)\cdot\frac{100\rho(1+r_F)-1}{(1+\rho)(1+r_F)}
                =\frac{\rho\bigl[100(1+r_F)+1\bigr]}{1+\rho}.
            \end{aligned}
        \]

        财富 \(w=100+\dfrac{1}{1+r_F}=\dfrac{100(1+r_F)+1}{1+r_F}\) 的分子正是上式的 \(100(1+r_F)+1\)，故
        \[
            c_0^{*}=\frac{w}{1+\rho},\qquad
            c_1^{*}=\frac{\rho(1+r_F)}{1+\rho}\,w.
        \]

        两式相除，零期与一期的最优消费只差偏好权重和总收益：
        \[
            c_1^{*}=\rho(1+r_F)\,c_0^{*}.
        \]

        因此最优现期消费占财富 $\dfrac{1}{1+\rho}$，这个份额不随 \(r_F\) 改变；最优 1 期消费占 \(\dfrac{\rho(1+r_F)}{1+\rho}\)，是在该份额上再乘以 \(\rho(1+r_F)\)。这是对数效用下的固定支出份额。最优效用
        \[
            U^{b}=\log\Bigl(\frac{w}{1+\rho}\Bigr)+\rho\log\Bigl(\frac{\rho(1+r_F)}{1+\rho}\,w\Bigr).
        \]

        \item \textbf{比较静态。}上面已经求出最优消费量：
        \[
            c_0^{*}=\frac{100(1+r_F)+1}{(1+\rho)(1+r_F)}
            =\frac{1}{1+\rho}\Bigl(100+\frac{1}{1+r_F}\Bigr),\qquad
            c_1^{*}=\frac{\rho\bigl[100(1+r_F)+1\bigr]}{1+\rho}.
        \]
        \(r_F\) 只出现在 \(c_0^{*}\) 的 \(1/(1+r_F)\) 里，故 \(r_F\) 上升时 \(c_0^{*}\) 下降；\(\rho\) 在分母，\(\rho\) 上升时 \(c_0^{*}\) 也下降。
        
        \(c_1^{*}\) 的分子含 \(\rho(1+r_F)\)，分母是 \(1+\rho\)，两项都使 \(c_1^{*}\) 随 \(r_F\)、\(\rho\) 上升。
        
        直觉：\(r_F\) 升高提高储蓄回报，\(\rho\) 升高加重未来消费权重，都会促使少消费于 0 期、多留于 1 期。
    \end{itemize}

    \item \textbf{$U^{b}\ge U^{a}$。}不能交易时只能消费禀赋，\((c_0,c_1)=(100,1)\)，效用为 \(U^{a}=\log 100\)。
    
    能交易时储蓄 \(S\) 可任取，其中 \(S=0\) 仍可行，对应的消费还是 \((100,1)\)，效用仍是 \(\log 100=U^{a}\)。\(U^{b}\) 是所有可行储蓄里的最大效用，而 \(S=0\) 只是其中一个可行点，最大值不会低于这一点：
    \[
        U^{b}=\max_{S\in\mathrm{R}}\,U\bigl(100-S,\,1+S(1+r_F)\bigr)
        \ge U(100,1)=U^{a}.
    \]
    等号当且仅当最优储蓄恰好是 \(S^{*}=0\)，即 \(100\rho(1+r_F)=1\)。

    \item \textbf{交易益处 $g$。}定义 $U^{b}(w-g)=U^{a}$：$w-g$ 是与交易后最优效用相同、但用一条 0 期财富数字表示的\textbf{确定性等价财富}；$g=w-(w-g)$ 即为因能交易而多出来的等价 0 期资源。

    将 $U^{a}=\log 100$ 与上式 $U^{b}$ 联立，记 $w-g=\hat w$，则
    \[
        (1+\rho)\log\hat w+\rho\log\bigl(\rho(1+r_F)\bigr)-(1+\rho)\log(1+\rho)=\log 100.
    \]

    取指数并整理：
    \[
        \hat w^{1+\rho}\,\frac{\bigl[\rho(1+r_F)\bigr]^{\rho}}{(1+\rho)^{1+\rho}}=100
        \quad\Rightarrow\quad
        \hat w=\rho^{-\frac{\rho}{1+\rho}}(1+r_F)^{-\frac{\rho}{1+\rho}}
        100^{\frac{1}{1+\rho}}(1+\rho).
    \]

    故
    \[
        g=w-\hat w
        =w-\rho^{-\frac{\rho}{1+\rho}}(1+r_F)^{-\frac{\rho}{1+\rho}}
        100^{\frac{1}{1+\rho}}(1+\rho).
    \]
    $g$ 随 $\rho$、$r_F$ 增加而增加：二者都提高把消费从 0 期转到 1 期」的动机，能在市场上做跨期配置带来的等价收益更大。
\end{enumerate}
\end{solution}


